% !PRINTED 717
\weberpage{717}

\begin{center}
\S.\,196.

\gesperrt{Classenzahl im Körper der achten Einheitswurzeln.}
\end{center}

Wir wollen, einerseits zur Veranschaulichung der bisherigen Resultate, andererseits um für die später anzuwendende vollständige Induction eine Basis zu gewinnen, den Fall $m = 8$, $\mu = 4$, $\nu = 2$ zunächst behandeln.

Hier haben wir ausser dem Hauptcharakter nur drei Charaktere, nämlich, wenn
\[
n \equiv (-1)^\alpha 5^\beta \pmod{8}
\]
ist,
\[
\begin{aligned}
\chi_1(n) & = (-1)^\beta \\
\chi_2(n) & = (-1)^\alpha \\
\chi_3(n) & = (-1)^{\alpha+\beta}.
\end{aligned}
\]
Für $n = -1$ ist $\alpha = 1$, $\beta = 0$, und folglich gehört $\chi_1$ zur ersten, $\chi_2$, $\chi_3$ zur zweiten Art [\S.\,195, (5)], und es ergeben sich für die vier Zahlen $n = 1, 3, 5, 7$ folgende Werthe dieser Charaktere:
\[
\begin{aligned}
\chi_1: & \quad +1, \; -1, \; -1, \; +1 \\
\chi_2: & \quad +1, \; -1, \; +1, \; -1 \\
\chi_3: & \quad +1, \; +1, \; -1, \; -1.
\end{aligned}
\]
Daraus erhält man die drei Functionen $f_1, f_2, f_3$ [\S.\,194, (2)]:
\[
\begin{aligned}
f_1(x) & = x - x^3 - x^5 + x^7 = x(1-x^2)(1-x^4) \\
f_2(x) & = x - x^3 + x^5 - x^7 = x(1-x^2)(1+x^4) \\
f_3(x) & = x + x^3 - x^5 - x^7 = x(1+x^2)(1-x^4).
\end{aligned}
\]
Für $r$ können wir setzen:
\[
r = \frac{1+i}{\sqrt{2}}, \quad r^2 = i, \quad r^3 = -\frac{1-i}{\sqrt{2}}, \quad r^4 = -1,
\]
und danach ergiebt sich:
\[
\begin{aligned}
f_1(r) & = 2\sqrt{2}, & f_2(r) & = 0, & f_3(r) & = 2\sqrt{2}\,i, \\
f_1(r^2) & = 0, & f_2(r^2) & = 4i, & f_3(r^2) & = 0, \\
f_1(r^3) & = -2\sqrt{2}, & f_2(r^3) & = 0, & f_3(r^3) & = 2\sqrt{2}\,i.
\end{aligned}
\]
Für $x = r^4$ verschwinden alle drei Functionen $f(x)$, und für $x = r^5, r^6, r^7$ erhält $f_1(x)$ dieselben, $f_2(x)$ und $f_3(x)$ die entgegengesetzten Werthe, wie für $x = r^3, r^2, r$.
